% ============================================================================
% 3.2 run.sh：交叉编译与启动
% ============================================================================
\subsection{为什么需要三步式构建}

在 macOS 上写一个在 QEMU 内跑的 ARM64 裸机程序，
你面对的是"跨平台"问题：
你的电脑是 ARM64 架构，但你写的程序跑在另一个"ARM64 机器"里。
实际上，它们\textbf{指令集是一样的}，区别只在地址空间和系统调用约定上。

具体来说：

\begin{itemize}
    \item macOS 的 AArch64 用户态程序通过 \texttt{svc \#0x80} 调用内核
    \item 裸机 AArch64 程序根本没有内核，也没有 svc 约定
    \item 因此我们需要一个"不链接任何系统库"的纯裸机二进制
\end{itemize}

因此构建是三步：

\begin{lstlisting}[style=bashstyle,caption=QEMU/run.sh — 编译脚本]
#!/bin/zsh
set -e
mkdir -p bin

# 第一步：用 clang 作为交叉汇编器
# --target=aarch64-none-elf 告诉 clang 目标是裸机 ARM64
#  -c 只编译不链接，产出 bin/start.o
/opt/homebrew/opt/llvm/bin/clang \
    --target=aarch64-none-elf \
    -c start.s \
    -o bin/start.o

# 第二步：用 LLVM 的 ld.lld 链接
# -Ttext=0x40000000 告诉链接器把 .text 段放到 0x4000\_0000
#    这是 QEMU virt 机器的内核默认加载地址
# -e _start 指定入口符号
/opt/homebrew/opt/lld/bin/ld.lld \
    -Ttext=0x40000000 \
    -e _start \
    -o bin/kernel.elf bin/start.o

# 第三步：启动 QEMU
# -machine virt ：使用 virt 虚拟机器（通用 ARM64 虚拟平台）
# -cpu cortex-a72 ：模拟 ARM Cortex-A72 处理器
# -nographic ：用当前终端作为串口控制台（不弹出图形窗口）
# -kernel bin/kernel.elf ：把 ELF 文件作为内核加载
qemu-system-aarch64 \
    -machine virt \
    -cpu cortex-a72 \
    -nographic \
    -kernel bin/kernel.elf
\end{lstlisting}

\paragraph{地址 0x4000\_0000 的来源}
QEMU 的 virt 机器把"内核可加载"的 RAM 起始地址固定为 \texttt{0x4000\_0000}
（1 GB 处）。你编译时让链接器把 \texttt{\_start} 放到这个地址，
然后 QEMU 从这个地址开始执行你的程序。

如果你想让程序执行在别的地址，可以随意改这个数字，
但必须保证它在 QEMU 的 DRAM 范围内（0x4000\_0000 到 0x8000\_0000 之间是一个安全的默认区）。

\paragraph{退出 QEMU}
在 \texttt{-nographic} 模式下，\textbf{Ctrl-A X}（先按 Ctrl-A，松开，再按 X）
可以退出 QEMU。Ctrl-C 是被送进模拟机器里的，不会让你退出 QEMU 主程序。

\subsection{常见错误}

\begin{itemize}
    \item \textbf{command not found: ld.lld}：你的 Homebrew 没装 lld。
        用 \texttt{brew install lld} 或直接 \texttt{brew install llvm} 装完整 LLVM。
    \item \textbf{No 'kernel' option found for this machine type}：QEMU 版本太老
        或你没装 AArch64 版的 QEMU。确认你装的是 \texttt{qemu-system-aarch64}。
    \item \textbf{程序跑了但没输出}：检查 UART 基地址是否为 \texttt{0x09000000}。
        不同机器类型使用不同地址（例如 \texttt{versatilepb} 用的是 \texttt{0x101f1000}）。
    \item \textbf{地址异常或崩溃}：链接地址与 QEMU 加载地址不一致。
        用 \texttt{-Ttext=0x40000000} 配合 \texttt{-kernel bin/kernel.elf}，
        会按链接地址原样加载各段到对应物理地址。
\end{itemize}

运行脚本后，终端会输出：

\begin{verbatim}
Hello QEMU AArch64 bare metal!
\end{verbatim}

然后程序进入 \texttt{wfe} 永久等待。按 Ctrl-A X 退出 QEMU。
